% =============================================================================
%  Elementos pós-textuais (NBR 14724): referências (NBR 6023:2018), glossário e
%  apêndices.
% =============================================================================

% ---------- Referências ----------
\titulopostextual{Referências}
\begingroup
\singlespacing
\setlength{\parindent}{0pt}
\setlength{\parskip}{12pt}
\raggedright

ALPHA CEPHEI. \textbf{Vosk speech recognition toolkit}. [\textit{S. l.}]: Alpha Cephei, 2026.
Disponível em: \url{https://alphacephei.com/vosk/}. Acesso em: 22 out. 2026.

ANTHROPIC. \textbf{Claude API documentation}. San Francisco: Anthropic, 2026. Disponível em:
\url{https://docs.anthropic.com/}. Acesso em: 22 out. 2026.

ASSOCIAÇÃO BRASILEIRA DE NORMAS TÉCNICAS. \textbf{NBR 6023}: informação e documentação: referências:
elaboração. Rio de Janeiro: ABNT, 2018.

ASSOCIAÇÃO BRASILEIRA DE NORMAS TÉCNICAS. \textbf{NBR 10520}: informação e documentação: citações em
documentos: apresentação. Rio de Janeiro: ABNT, 2023.

ASSOCIAÇÃO BRASILEIRA DE NORMAS TÉCNICAS. \textbf{NBR 14724}: informação e documentação: trabalhos
acadêmicos: apresentação. Rio de Janeiro: ABNT, 2011.

BECK, Kent. \textbf{Test-driven development}: by example. Boston: Addison-Wesley, 2003.

COCKBURN, Alistair. \textbf{Hexagonal architecture}. [\textit{S. l.}], 2005. Disponível em:
\url{https://alistair.cockburn.us/hexagonal-architecture/}. Acesso em: 22 out. 2026.

COHN, Mike. \textbf{Succeeding with agile}: software development using Scrum. Upper Saddle River:
Addison-Wesley, 2009.

DEBIAN. \textbf{Debian policy manual}. [\textit{S. l.}]: The Debian Project, 2026. Disponível em:
\url{https://www.debian.org/doc/debian-policy/}. Acesso em: 22 out. 2026.

EVANS, Eric. \textbf{Domain-driven design}: tackling complexity in the heart of software. Boston:
Addison-Wesley, 2004.

FORD, Neal; PARSONS, Rebecca; KUA, Patrick. \textbf{Building evolutionary architectures}: support
constant change. Sebastopol: O'Reilly Media, 2017.

FOWLER, Martin. \textbf{Inversion of control containers and the dependency injection pattern}.
[\textit{S. l.}], 2004. Disponível em: \url{https://martinfowler.com/articles/injection.html}. Acesso
em: 22 out. 2026.

FOWLER, Martin. \textbf{Refactoring}: improving the design of existing code. 2. ed. Boston:
Addison-Wesley, 2018.

FREEDESKTOP.ORG. \textbf{Desktop entry specification}: version 1.5. [\textit{S. l.}], 2020.
Disponível em: \url{https://specifications.freedesktop.org/desktop-entry-spec/latest/}. Acesso em:
22 out. 2026.

GAMMA, Erich \textit{et al.} \textbf{Design patterns}: elements of reusable object-oriented software.
Reading: Addison-Wesley, 1994.

HARRIS, Charles R. \textit{et al.} Array programming with NumPy. \textbf{Nature}, London, v. 585,
p. 357-362, 2020.

HUNT, Andrew; THOMAS, David. \textbf{The pragmatic programmer}: from journeyman to master. Reading:
Addison-Wesley, 1999.

INTERNATIONAL ORGANIZATION FOR STANDARDIZATION. \textbf{ISO/IEC 25010:2011}: systems and software
engineering: systems and software quality requirements and evaluation (SQuaRE): system and software
quality models. Geneva: ISO, 2011.

JURAFSKY, Daniel; MARTIN, James H. \textbf{Speech and language processing}: an introduction to natural
language processing, computational linguistics, and speech recognition. 2. ed. Upper Saddle River:
Prentice Hall, 2009.

LARMAN, Craig; BASILI, Victor R. Iterative and incremental development: a brief history.
\textbf{Computer}, Los Alamitos, v. 36, n. 6, p. 47-56, 2003.

LEVKIVSKYI, Ivan; LEHTOSALO, Jukka; LANGA, Łukasz. \textbf{PEP 544}: protocols: structural subtyping
(static duck typing). [\textit{S. l.}]: Python Software Foundation, 2017. Disponível em:
\url{https://peps.python.org/pep-0544/}. Acesso em: 22 out. 2026.

LISKOV, Barbara H.; WING, Jeannette M. A behavioral notion of subtyping. \textbf{ACM Transactions on
Programming Languages and Systems}, New York, v. 16, n. 6, p. 1811-1841, 1994.

MARTIN, Robert C. \textbf{Agile software development}: principles, patterns, and practices. Upper
Saddle River: Prentice Hall, 2003.

MARTIN, Robert C. \textbf{Código limpo}: habilidades práticas do agile software. Rio de Janeiro: Alta
Books, 2009.

MARTIN, Robert C. \textbf{Arquitetura limpa}: o guia do artesão para estrutura e design de software.
Rio de Janeiro: Alta Books, 2019.

McCABE, Thomas J. A complexity measure. \textbf{IEEE Transactions on Software Engineering}, New York,
v. SE-2, n. 4, p. 308-320, 1976.

MESZAROS, Gerard. \textbf{xUnit test patterns}: refactoring test code. Upper Saddle River:
Addison-Wesley, 2007.

MEYER, Bertrand. \textbf{Object-oriented software construction}. Englewood Cliffs: Prentice Hall,
1988.

NATIONAL INSTITUTE OF STANDARDS AND TECHNOLOGY. \textbf{FIPS PUB 180-4}: secure hash standard (SHS).
Gaithersburg: NIST, 2015.

NIELSEN, Jakob. \textbf{Usability engineering}. Boston: Academic Press, 1993.

OLLAMA. \textbf{Ollama API}. [\textit{S. l.}], 2026. Disponível em:
\url{https://github.com/ollama/ollama/blob/main/docs/api.md}. Acesso em: 26 out. 2026.

PIPEWIRE. \textbf{PipeWire documentation}. [\textit{S. l.}], 2026. Disponível em:
\url{https://docs.pipewire.org/}. Acesso em: 26 out. 2026.

POVEY, Daniel \textit{et al.} The Kaldi speech recognition toolkit. \textit{In}: IEEE WORKSHOP ON
AUTOMATIC SPEECH RECOGNITION AND UNDERSTANDING, 2011, Waikoloa. \textbf{Proceedings} [...].
[\textit{S. l.}]: IEEE Signal Processing Society, 2011.

PRESSMAN, Roger S.; MAXIM, Bruce R. \textbf{Engenharia de software}: uma abordagem profissional.
8. ed. Porto Alegre: AMGH, 2016.

PRESTON-WERNER, Tom. \textbf{Semantic versioning 2.0.0}. [\textit{S. l.}], 2013. Disponível em:
\url{https://semver.org/}. Acesso em: 22 out. 2026.

SOHN, Jongseo; KIM, Nam Soo; SUNG, Wonyong. A statistical model-based voice activity detection.
\textbf{IEEE Signal Processing Letters}, New York, v. 6, n. 1, p. 1-3, 1999.

SOMMERVILLE, Ian. \textbf{Engenharia de software}. 10. ed. São Paulo: Pearson Education do Brasil,
2018.

SQLITE. \textbf{SQLite documentation}. [\textit{S. l.}], 2026. Disponível em:
\url{https://www.sqlite.org/docs.html}. Acesso em: 22 out. 2026.

VASWANI, Ashish \textit{et al.} Attention is all you need. \textit{In}: CONFERENCE ON NEURAL
INFORMATION PROCESSING SYSTEMS, 31., 2017, Long Beach. \textbf{Advances in neural information
processing systems 30}. Red Hook: Curran Associates, 2017. p. 5998-6008.

\endgroup

% ---------- Glossário ----------
\titulopostextual{Glossário}
\begingroup
\singlespacing
\setlength{\parindent}{0pt}
\setlength{\parskip}{8pt}

\textbf{Adaptador}: componente que implementa uma porta usando uma tecnologia concreta, como um banco de
dados ou um serviço de voz.

\textbf{Avatar}: representação gráfica do assistente; neste trabalho, um busto gerado a partir de
fotografias do usuário.

\textbf{Cache LRU}: armazenamento temporário que, quando cheio, descarta primeiro os itens usados há
mais tempo.

\textbf{Cadeia de resumos}: sequência de registros em que cada um guarda o resumo criptográfico do
anterior, tornando detectável qualquer alteração.

\textbf{Dublê de teste}: objeto que substitui um colaborador real durante um teste, como um sintetizador
que apenas registra os textos recebidos.

\textbf{Fluxo contínuo} (\en{streaming}): entrega de um resultado em partes, à medida que é produzido.

\textbf{Frase de ativação}: expressão que chama o assistente sem uso das mãos; neste trabalho, “ok Fael”.

\textbf{Frase distratora}: frase parecida com a de ativação, incluída na gramática para absorver o que
não deve ativar o assistente.

\textbf{Função de aptidão arquitetural}: verificação objetiva, preferencialmente automatizada, de uma
característica desejada da arquitetura.

\textbf{Gramática restrita}: lista fechada de frases entre as quais o reconhecedor de fala deve escolher.

\textbf{Pipeline}: encadeamento de estágios que processam dados em sequência, de forma sobreposta no
tempo.

\textbf{Porta}: contrato que descreve o que a aplicação precisa do mundo externo, ou o que oferece a ele.

\textbf{Raiz de composição}: ponto único do programa em que as implementações concretas são criadas e
conectadas.

\textbf{Token de cancelamento}: sinal compartilhado pelos estágios de uma interação, cuja ativação
interrompe todos eles.

\textbf{\en{Wheel}}: formato de distribuição de pacotes Python já construídos, instaláveis sem
compilação.

\endgroup

% ---------- Apêndices ----------
\clearpage
\phantomsection
\renewcommand{\thequadro}{A.\arabic{quadro}}
\setcounter{quadro}{0}
\titulopostextual{Apêndice A – Rastreabilidade entre requisitos e testes}

Este apêndice associa cada requisito à suíte que o verifica. Os nomes das suítes correspondem aos
arquivos da pasta \cod{tests}.

\begin{singlespace}
\legendaquadro{Rastreabilidade dos requisitos funcionais}
\small
\begin{tabularx}{\textwidth}{|P{2.4cm}|L|}
\hline
\textbf{Requisitos} & \textbf{Suítes de teste} \\ \hline
\endfirsthead
\hline
\textbf{Requisitos} & \textbf{Suítes de teste} \\ \hline
\endhead
RF01, RF14 & \cod{test_avatar}, \cod{test_configuracoes} \\ \hline
RF02 & \cod{test_animacao}, \cod{test_analise_audio} \\ \hline
RF03, RF09, RF13 & \cod{test_orquestrador}, \cod{test_vad} \\ \hline
RF04, RF05 & \cod{test_ativacao}, \cod{test_ganho}, \cod{test_sessao_pipeline_vigia} \\ \hline
RF06 & \cod{test_adaptadores_voz}, \cod{test_perfil_voz}, \cod{test_configuracoes} \\ \hline
RF07, RF08, RF10, RF12 & \cod{test_bootstrap_ui}, \cod{test_orquestrador}, \cod{test_estados_barramento}, \cod{test_configuracoes} \\ \hline
RF11 & \cod{test_bootstrap_ui}, \cod{test_i18n} \\ \hline
RF15 a RF17, RF19 & \cod{test_agenda}, \cod{test_modelo_offline}, \cod{test_ferramentas} \\ \hline
RF18 & \cod{test_agendador}, \cod{test_modelo_offline} \\ \hline
RF20, RF22, RF23 & \cod{test_sqlite}, \cod{test_repositorio_json} \\ \hline
RF21 & \cod{test_memoria}, \cod{test_i18n} \\ \hline
RF24, RF26 & \cod{test_modelo_claude}, \cod{test_ferramentas} \\ \hline
RF25, RF27 & \cod{test_modelo_offline}, \cod{test_bootstrap_ui} \\ \hline
RF28, RF29 & \cod{test_programas} \\ \hline
RF30 & \cod{test_perfil_voz}, \cod{test_analise_audio} \\ \hline
RF31, RF32 & \cod{test_configuracoes}, \cod{test_sessao_pipeline_vigia} \\ \hline
RF33, RF34 & \cod{test_i18n} \\ \hline
RF35 & \cod{test_configuracao_tema}, \cod{test_i18n} \\ \hline
RF36 & roteiro \cod{empacotamento/testar_deb.sh} (\cod{make testar-deb}) \\ \hline
RF37 & \cod{test_ollama} \\ \hline
RF38 & \cod{test_banco_e_modelos}, \cod{test_configuracoes} \\ \hline
RF39, RF40 & \cod{test_microfone_inicializacao}, \cod{test_configuracoes} \\ \hline
RF41, RF42 & \cod{test_banco_e_modelos}, \cod{test_configuracoes} \\ \hline
\end{tabularx}
\fonte{elaboração própria (2026).}
\end{singlespace}

\begin{singlespace}
\legendaquadro{Rastreabilidade dos requisitos não funcionais}
\small
\begin{tabularx}{\textwidth}{|P{2.4cm}|L|}
\hline
\textbf{Requisitos} & \textbf{Verificação} \\ \hline
\endfirsthead
\hline
\textbf{Requisitos} & \textbf{Verificação} \\ \hline
\endhead
RNF01 & \cod{test_bootstrap_ui} e medição de tempo \\ \hline
RNF02, RNF03 & \cod{test_orquestrador}, \cod{test_sessao_pipeline_vigia} \\ \hline
RNF04 & \cod{test_agendador} \\ \hline
RNF05 & \cod{test_animacao}, \cod{test_avatar}, \cod{test_adaptadores_voz} \\ \hline
RNF06 & \cod{test_sqlite} \\ \hline
RNF07 & \cod{test_modelo_claude} \\ \hline
RNF08 & \cod{test_bootstrap_ui}, \cod{test_ativacao}, \cod{test_adaptadores_voz} \\ \hline
RNF09 & \cod{test_orquestrador}, \cod{test_estados_barramento} \\ \hline
RNF10, RNF12, RNF17 & revisão de código e do catálogo de mensagens \\ \hline
RNF11 & \cod{test_ferramentas}, \cod{test_programas} \\ \hline
RNF13 & \cod{test_ativacao} \\ \hline
RNF14 a RNF16 & \cod{test_configuracoes}, \cod{test_bootstrap_ui} \\ \hline
RNF18 & \cod{test_i18n} \\ \hline
RNF19, RNF22 & \cod{test_arquitetura} \\ \hline
RNF20, RNF21 & ruff, mypy e pytest-cov em \cod{make verificar} e na integração contínua \\ \hline
RNF23, RNF24 & integração contínua (Python 3.10 e 3.12) e \cod{make testar-deb} \\ \hline
RNF25 & revisão do arquivo \cod{CHANGELOG.md} \\ \hline
RNF26, RNF27 & \cod{test_renderizacao}, \cod{test_analise_audio}, \cod{test_adaptadores_voz} e medição no app real \\ \hline
RNF28 & \cod{test_banco_e_modelos} e a fixture \cod{casa_isolada} (\cod{tests/conftest.py}) \\ \hline
\end{tabularx}
\fonte{elaboração própria (2026).}
\end{singlespace}

\clearpage
\renewcommand{\thequadro}{B.\arabic{quadro}}
\setcounter{quadro}{0}
\titulopostextual{Apêndice B – Comandos de voz}

O Quadro~B.1 apresenta exemplos de comandos aceitos. Os comandos são entendidos nos dois idiomas,
independentemente do idioma das respostas. Com o modelo de linguagem ativo, qualquer formulação em
linguagem natural é aceita; os exemplos abaixo também funcionam no modo offline, exceto quando indicado.

\begin{singlespace}
\legendaquadro{Exemplos de comandos de voz}
\small
\begin{tabularx}{\textwidth}{|P{3.4cm}|L|L|}
\hline
\textbf{Função} & \textbf{Português} & \textbf{Inglês} \\ \hline
\endfirsthead
\hline
\textbf{Função} & \textbf{Português} & \textbf{Inglês} \\ \hline
\endhead
Chamar o assistente & ok Fael & ok Fael \\ \hline
Chamar e comandar & ok Fael, anota tarefa comprar pão & ok Fael, add task buy bread \\ \hline
Criar tarefa & anota tarefa revisar o relatório & add task review the report \\ \hline
Listar tarefas & quais são minhas tarefas? & what are my tasks? \\ \hline
Concluir tarefa & terminei a tarefa 2 & complete task 2 \\ \hline
Lembrete relativo & me lembra daqui a 20 minutos de beber água & remind me in 20 minutes to drink water \\ \hline
Lembrete em data e hora (requer o modelo) & amanhã às 9 me lembra da reunião & tomorrow at 9 remind me of the meeting \\ \hline
Histórico & o que eu fiz ontem? o que concluí na semana passada? & what did I do yesterday? what did I do last week? \\ \hline
Abrir programa & abre o navegador; abrir a calculadora & open the browser; open the calculator \\ \hline
Hora e data & que horas são? que dia é hoje? & what time is it? what's the date? \\ \hline
Nova conversa & nova conversa & new conversation \\ \hline
Encerrar o modo contínuo & tchau & bye \\ \hline
\end{tabularx}
\fonte{elaboração própria (2026).}
\end{singlespace}

\clearpage
\renewcommand{\thequadro}{C.\arabic{quadro}}
\setcounter{quadro}{0}
\titulopostextual{Apêndice C – Configuração}

O Quadro~C.1 relaciona as variáveis de ambiente e as opções de linha de comando. A precedência é:
preferência salva na tela de configurações, depois opção de linha de comando ou variável de ambiente e,
por fim, valor padrão.

\begin{singlespace}
\legendaquadro{Variáveis de ambiente e opções de linha de comando}
\small
\begin{tabularx}{\textwidth}{|P{4.6cm}|P{3.2cm}|L|}
\hline
\textbf{Variável ou opção} & \textbf{Padrão} & \textbf{Efeito} \\ \hline
\endfirsthead
\hline
\textbf{Variável ou opção} & \textbf{Padrão} & \textbf{Efeito} \\ \hline
\endhead
\cod{ANTHROPIC_API_KEY} & -- & Ativa o Claude; sem ela, modelo local (Ollama) ou modo offline \\ \hline
\cod{ROBOFAEL_CEREBRO} & auto & \cod{auto} (Claude, depois Ollama, depois offline), \cod{claude}, \cod{ollama} ou \cod{offline} \\ \hline
\cod{ROBOFAEL_OLLAMA_URL} & \cod{http://localhost:11434} & Servidor do Ollama \\ \hline
\cod{ROBOFAEL_OLLAMA_MODELO} & \cod{llama3.2:3b} & Modelo local (sem ele, o primeiro instalado que aceite ferramentas) \\ \hline
\cod{ROBOFAEL_USUARIO} & Rafael & Nome usado na saudação e nas instruções do modelo \\ \hline
\cod{ROBOFAEL_ASSISTENTE} & Robofael & Nome do assistente \\ \hline
\cod{ROBOFAEL_FRASE_ATIVACAO} & ok Fael & Frase de ativação \\ \hline
\cod{ROBOFAEL_ATIVACAO} ou \cod{--sem-ativacao} & 1 (ligada) & Liga ou desliga a escuta passiva \\ \hline
\cod{ROBOFAEL_IDIOMA} ou \cod{--idioma} & pt-BR & Idioma das respostas e da interface (pt-BR ou en-US) \\ \hline
\cod{ROBOFAEL_MODELO} & \cod{claude-opus-5-5} & Modelo do Claude \\ \hline
\cod{ROBOFAEL_VOZ} & \cod{pt-BR-AntonioNeural} & Voz neural padrão \\ \hline
\cod{ROBOFAEL_DADOS} & \cod{./dados} (instalado: \cod{~/.local/share/robofael}) & Banco de dados (que pode ser movido nas configurações), perfil e gravações da voz \\ \hline
\cod{ROBOFAEL_CACHE} & \cod{./.cache} (instalado: \cod{~/.cache/robofael}) & Quadros do avatar e falas sintetizadas \\ \hline
\cod{ROBOFAEL_ASSETS} & \cod{./assets} & Fotografias dos avatares \\ \hline
\cod{ROBOFAEL_MODELOS} & \cod{./modelos} & Modelo de reconhecimento local \\ \hline
\cod{--avatar} & o salvo & Avatar inicial (crachá ou \en{selfie}) \\ \hline
\cod{--continuo} & desligado & Volta a ouvir após cada resposta \\ \hline
\cod{-v} & desligado & Registros detalhados (texto ouvido, pico e ganho) \\ \hline
\end{tabularx}
\fonte{elaboração própria (2026).}
{\noindent\small\singlespacing Nota: os nomes anteriores à versão 3.5.0 (\cod{EUROBO_*}) ainda são aceitos,
com aviso de depreciação, quando a variável \cod{ROBOFAEL_*} correspondente não está definida.\par}
\end{singlespace}

\clearpage
\titulopostextual{Apêndice D – Procedimentos de instalação, execução e verificação}

\textbf{Instalação pelo pacote (Ubuntu 24.04).} Na pasta do projeto, executar \cod{make deb} para gerar
o pacote em \cod{dist/}. Em seguida, instalá-lo com
\cod{sudo apt install ./dist/robofael_3.8.0_amd64.deb}. O gerenciador de pacotes instala as dependências
de sistema, e o pacote cria o ambiente Python em \cod{/opt/robofael} sem acessar a internet. O
assistente é aberto pelo item Robofael do menu de aplicativos ou pelo comando \cod{robofael}; a
ferramenta de gravação da voz, pelo item correspondente ou pelo comando \cod{robofael-gravar-voz}. Para
remover, usar \cod{sudo apt remove robofael}, que preserva os dados, ou \cod{sudo apt purge robofael},
que também apaga a pasta do programa; os dados pessoais só são removidos manualmente.

\textbf{Configurações de uso.} O botão de engrenagem abre a tela de configurações. Além de idioma,
avatar, voz e ativação por voz, ela permite escolher o microfone; marcar as opções de abrir o assistente
ao entrar no computador e de usar o comando \cod{robofael} em qualquer terminal (no pacote, o comando já
existe); ver o modelo de linguagem em uso e os modelos compatíveis do Claude e do Ollama; e, na seção de
memória, mudar o banco de pasta ou fazer um backup, operações que valem na hora. Para restaurar um
backup, basta fechar o assistente e copiar o arquivo do backup sobre o \cod{memoria.db}. Como o
assistente aberto pelo login não lê o \cod{~/.bashrc}, a chave do Claude, nesse caso, deve ser definida
em \cod{~/.config/environment.d/}.

\textbf{Execução em modo de desenvolvimento.} Executar \cod{make instalar} para criar o ambiente com as
dependências de desenvolvimento, \cod{make baixar-modelo-voz} para obter o modelo de reconhecimento local
e \cod{make executar} para abrir o assistente. Opções podem ser passadas com
\cod{make executar ARGS="--idioma en-US -v"}. A gravação da voz é feita com \cod{make gravar-voz}.

\textbf{Verificação.} O comando \cod{make verificar} executa a análise estática, a verificação de tipos
e os 643 testes; \cod{make testar-rapido} omite os testes de interface e de renderização;
\cod{make cobertura} gera o relatório de cobertura; \cod{make integridade} verifica a cadeia de resumos
do histórico; e \cod{make testar-deb} constrói o pacote e o testa em uma pasta temporária.

\textbf{Geração deste documento.} O texto desta monografia está em \cod{docs/monografia}, em LaTeX, e é
compilado com \cod{make monografia}, que produz o arquivo \cod{docs/monografia/monografia.pdf}.
